En programación competitiva, dominar o al menos tener buenas nociones de los conceptos principales de teoría de números no es un lujo, sino una necesidad. A lo largo de este texto y en la práctica dentro de la ICPC, el lector encontrará problemas donde aparecen de manera natural ideas como divisibilidad, números primos, residuos, máximo común divisor y aritmética modular. En términos generales, la teoría de números estudia las propiedades de los enteros, y por ello ofrece un lenguaje muy útil para describir, simplificar y resolver muchos problemas clásicos de concurso.
\section{Bases}
Decimos que $a$ \textbf{divide} a $b$ (escrito $a \mid b$) si existe un entero $k$ tal que $b = ka$. Por ejemplo, $4 \mid 20$ porque $20 = 4 \cdot 5$, de modo que aquí el entero que funciona es $k = 5$. En cambio, $6 \nmid 20$ porque no existe un entero $k$ tal que $20 = 6k$.

Veamos ahora algunas propiedades de la divisibilidad que aparecen con frecuencia en competencia.

\textbf{Propiedad 1.} Si un mismo número divide a otros dos, entonces también divide su suma y su diferencia.
\[
a \mid b \text{ y } a \mid c \implies a \mid (b \pm c)
\]

Por ejemplo, como $3 \mid 12$ y $3 \mid 21$, entonces también $3 \mid (21-12) = 9$ y $3 \mid (21+12) = 33$.

\textbf{Propiedad 2.} Si un número divide a otro, entonces también divide cualquier múltiplo de ese segundo número.
\[
a \mid b \implies a \mid bc
\]

Por ejemplo, como $4 \mid 12$, también se cumple que $4 \mid (12 \cdot 5) = 60$.

\textbf{Propiedad 3.} Si un número divide a un segundo y ese segundo divide a un tercero, entonces el primero divide al tercero.
\[
a \mid b \text{ y } b \mid c \implies a \mid c
\]

Por ejemplo, como $2 \mid 8$ y $8 \mid 40$, se sigue que $2 \mid 40$.


El \textbf{máximo común divisor} de dos números $a$ y $b$, $\gcd(a,b)$ por sus siglas en inglés, es el mayor entero que divide a ambos. El \textbf{mínimo común múltiplo} $\text{lcm}(a,b)$ es el menor entero positivo divisible por ambos.

Por ejemplo, si tomamos $a = 12$ y $b = 18$, entonces los divisores de $12$ son
\[
\{1,2,3,4,6,12\}
\]
y los divisores de $18$ son
\[
\{1,2,3,6,9,18\}.
\]
Los divisores comunes son $\{1,2,3,6\}$, así que $\gcd(12,18) = 6$.

De manera análoga, los primeros múltiplos positivos de $12$ son
\[
\{12,24,36,48,\dots\}
\]
y los de $18$ son
\[
\{18,36,54,\dots\}.
\]
El primero que aparece en ambas listas es $36$, por lo que $\text{lcm}(12,18) = 36$.

Podemos registrar además la siguiente relación importante entre ambas cantidades.

\textbf{Propiedad 4.} Para enteros positivos $a$ y $b$,
\[
\text{lcm}(a,b) = \frac{a \cdot b}{\gcd(a,b)}
\]

En el ejemplo anterior se obtiene
\[
\text{lcm}(12,18) = \frac{12 \cdot 18}{\gcd(12,18)} = \frac{216}{6} = 36.
\]

Podría parecer que esto nos da una implementación directa de $\text{lcm}$, pero aquí aparece una diferencia importante entre la teoría y la práctica. Matemáticamente,
\[
\text{lcm}(a,b) = \frac{a \cdot b}{\gcd(a,b)} = \left(\frac{a}{\gcd(a,b)}\right) \cdot b
\]
son expresiones equivalentes. Sin embargo, en una implementación real el orden de las operaciones sí importa. Si se usa un \texttt{int} de $32$ bits, normalmente se pueden representar valores entre $-2^{31}$ y $2^{31}-1$, es decir, aproximadamente hasta $10^9$ en magnitud. Con \texttt{long long}, que usualmente ocupa $64$ bits, el rango llega de $-2^{63}$ a $2^{63}-1$, aproximadamente hasta $10^{18}$.

Por ejemplo, si $a = 10^{10}$ y $b = 10^9$, entonces el producto $ab = 10^{19}$ ya no cabe en un \texttt{long long}, y si intentáramos hacer esa multiplicación primero se produciría un \textit{overflow}.\footnote{En programación, \textit{overflow} significa que el resultado de una operación queda fuera del rango que el tipo de dato puede representar. Cuando eso ocurre, el valor almacenado deja de ser confiable.} Así, aunque la fórmula sea correcta en el papel, calcular primero $a \cdot b$ puede producir un error antes de que siquiera se haga la división. Por esta razón, la forma recomendada de implementarlo es
\[
\boxed{\texttt{lcm(a, b) = a / gcd(a, b) * b}}
\]
porque primero se divide y así se reduce el tamaño de los valores intermedios.

Hasta este punto ya tenemos una relación entre $\text{lcm}(a,b)$ y $\gcd(a,b)$, pero todavía no hemos explicado cómo calcular alguna de estas dos cantidades. Empezaremos con el algoritmo de Euclides, que justamente permite obtener $\gcd(a,b)$ de manera eficiente.

\[
\gcd(a, b) = \gcd(b, a \bmod b)
\]
La idea es que cualquier divisor común de $a$ y $b$ también divide a $a \bmod b = a - \lfloor a/b \rfloor \cdot b$, y viceversa. Por ejemplo, si $a = 30$ y $b = 18$, entonces
\[
30 \bmod 18 = 12 = 30 - 1 \cdot 18.
\]
Los divisores comunes de $30$ y $18$ son $\{1,2,3,6\}$, y esos mismos números también dividen a $12$. De hecho,
\[
\gcd(30,18) = \gcd(18,12) = 6.
\]
Se repite este proceso hasta que el residuo es $0$, momento en que el divisor es el último valor no nulo. Esto nos deja una implementación muy corta y directa:
 
\begin{lstlisting}[language=C++]
long long gcd(long long a, long long b) {
    return b == 0 ? a : gcd(b, a % b);
}
 
long long lcm(long long a, long long b) {
    return a / gcd(a, b) * b;
}
\end{lstlisting}
 
La complejidad es $O(\log \min(a,b))$. En C++17 puedes usar directamente \texttt{gcd}.
\section{Aritmética modular}
Iniciemos analizando un ejemplo.
\\\\
\textbf{Ejemplos}
\textbf{[3.1] }Supongamos que hoy es Lunes, ¿qué día de la semana será dentro de 300 días?
\\
\textbf{Solución.} Podríamos hacer una lista de 300 números después del lunes, asignando cada día a cada número y ver cuál queda en el número 300, pero esto sería muy tardado y realmente innecesario, pues si notamos que cada 7 días se repite el mismo día, podemos ver que ya que al dividir 300 entre 7 nos da 42 y nos queda de residuo 6, por lo que sabemos que el día 294, es decir 42 veces 7, será Lunes y sólo tenemos que ver los restantes 6 días, que sabiendo que 6 días después de Lunes es Domingo, tendremos nuestra respuesta.
\\\\\
En este ejemplo vemos como la noción de que conjuntos de números se conformen de forma cíclica nos simplifica problemas. Si en el ejemplo anterior en lugar 300 días fueran 314 días, la respuesta sería la misma, podría parecer una afirmación fuerte, pero realmente al dividir 314 entre 7 como hicimos en el ejemplo notaremos que el residuo es el mismo, es decir 6, entonces el día 308, es decir 44 veces 7, será Lunes y tenemos nuevamente que ver los siguientes 6 días, esto nos lleva a pensar que si de días de la semana estamos hablando, dos días representan el mismo día de la semana si y sólo si tienen el mismo residuo al dividirlos entre 7.
\\Resulta que esto mismo lo podemos hacer con cualquier número natural $n$. Si $x$ y $y$ son dos números enteros tales que tienen mismo residuo al dividirles entre $n$, diremos que $x$ es congruente a $y$ módulo $n$, lo que podemos escribir con la siguiente notación:
\[x \equiv y \quad mod \quad n\]
Por ejemplo, $17 \equiv 5 \pmod{12}$ porque al dividir $17$ entre $12$ el residuo es $5$, y al dividir $5$ entre $12$ el residuo también es $5$. De manera similar, $31 \equiv 3 \pmod{7}$ porque ambos dejan residuo $3$ al dividirse entre $7$.

Veamos ahora algunas propiedades de las congruencias. Si $a\equiv b \, mod \, n$ y $c\equiv d \, mod \, n$ y $m$ es un entero positivo, entonces lo siguiente se cumple:

\textbf{Propiedad 1.} La congruencia respeta la suma.
\[
a+c \equiv b+d \, mod \, n
\]
Por ejemplo, como $17 \equiv 5 \pmod{12}$ y $8 \equiv 20 \pmod{12}$, se sigue que
\[
17 + 8 \equiv 5 + 20 \pmod{12}.
\]

\textbf{Propiedad 2.} La congruencia respeta el producto.
\[
ac \equiv bd \, mod \, n
\]
Por ejemplo, usando otra vez $17 \equiv 5 \pmod{12}$ y $8 \equiv 20 \pmod{12}$,
\[
17 \cdot 8 \equiv 5 \cdot 20 \pmod{12}.
\]

\textbf{Propiedad 3.} La congruencia respeta las potencias.
\[
a^m \equiv b^m \, mod \, n
\]
Por ejemplo, como $10 \equiv 3 \pmod{7}$, entonces
\[
10^2 \equiv 3^2 \pmod{7}.
\]
Como podemos notar, las congruencias respetan las sumas, productos y potencias como estamos acostumbrados, y en realidad para sumas y productos ocurre que cualquier número puede sustituirse por otro al que sea congruente en el módulo de la congruencia.
\\\\
Nótese que las congruencias son especialmente útiles cuando queremos evitar \textit{overflow}. Si estamos trabajando módulo $m$, podemos ir reduciendo cada resultado intermedio a su residuo módulo $m$ para que los números no crezcan más de lo necesario. En C++, dado un valor $a$, esto suele hacerse escribiendo \texttt{a \% m}.

Una aplicación fundamental de esta idea es la \textbf{exponenciación modular}, es decir, el problema de calcular $a^b \bmod m$ de manera eficiente. Si quisiéramos obtener, por ejemplo, $2^7$, una forma directa sería hacer
\[
2 \cdot 2 \cdot 2 \cdot 2 \cdot 2 \cdot 2 \cdot 2,
\]
lo cual requiere $6$ multiplicaciones. En cambio, si reutilizamos resultados intermedios, podemos hacer
\[
2^2 = 4, \qquad 2^4 = (2^2)^2 = 16, \qquad 2^7 = 2^4 \cdot 2^2 \cdot 2,
\]
y con ello reducimos notablemente el número de operaciones. Esta es la idea detrás de la exponenciación rápida: en lugar de construir la potencia multiplicando uno por uno, aprovechamos que podemos ir partiendo el exponente y recombinando resultados ya calculados.

La identidad clave es:
\[
a^b = \begin{cases} (a^{b/2})^2 & \text{si } b \text{ es par} \\ a \cdot a^{b-1} & \text{si } b \text{ es impar} \end{cases}
\]
Gracias a esto se obtiene un algoritmo de complejidad $O(\log b)$. Además, en cada paso podemos tomar módulo $m$, lo que da lugar justamente a la exponenciación modular:
 
\begin{lstlisting}[language=C++]
long long power(long long a, long long b, long long mod) {
    long long res = 1;
    a %= mod;
    while (b > 0) {
        //si b es impar se multiplica una vez extra
        if (b % 2) res = res * a % mod;
        a = a * a % mod;
        b = b/2;
    }
    return res;
}
\end{lstlisting}

\section{Números primos y criba de Eratóstenes}
 
Un número $p > 1$ es \textbf{primo} si sus únicos divisores son $1$ y $p$. Por ejemplo, $2$, $3$, $5$, $7$ y $11$ son números primos, mientras que $4$ no lo es porque puede dividirse entre $1$, $2$ y $4$.

\textbf{Teorema Fundamental de la Aritmética.} Todo entero $n > 1$ se puede escribir de forma única como producto de primos.

Por ejemplo,
\[
12 = 2^2 \cdot 3
\]
y
\[
60 = 2^2 \cdot 3 \cdot 5.
\]
La unicidad significa que, salvo el orden de los factores, no existe otra forma de escribir estos números como producto de primos.
 
Para verificar si $n$ es primo basta revisar divisores hasta $\sqrt{n}$. La razón es la siguiente: si $n$ no fuera primo, entonces podría escribirse como $n = a \cdot b$ con $1 < a < n$ y $1 < b < n$. Si ocurriera que tanto $a > \sqrt{n}$ como $b > \sqrt{n}$, entonces
\[
a \cdot b > \sqrt{n}\cdot \sqrt{n} = n,
\]
lo cual es imposible. Por lo tanto, todo número compuesto tiene al menos un divisor menor o igual que $\sqrt{n}$.

Por ejemplo, si $n = 91$, entonces
\[
91 = 7 \cdot 13.
\]
Aquí uno de los divisores, $7$, ya es menor que $\sqrt{91}$, así que basta revisar hasta ese rango para detectar que $91$ no es primo. En consecuencia, si ningún entero entre $2$ y $\lfloor \sqrt{n} \rfloor$ divide a $n$, entonces $n$ es primo. Esto da un algoritmo de complejidad $O(\sqrt{n})$.
 
\begin{lstlisting}[language=C++]
bool prime(long long n) {
    if (n < 2) return false;
    if (n % 2 == 0) return n == 2;
    for (long long i = 3; i * i <= n; i += 2) {
        if (n % i == 0) return false;
    }
    return true;
}
\end{lstlisting}
 
Para encontrar todos los primos hasta $n$, la \textbf{criba de Eratóstenes} es el método estándar. La idea central es mantener una tabla de números y tachar los múltiplos de cada primo que vayamos encontrando. Se comienza con el $2$, se reconocen como compuestos sus múltiplos $4,6,8,\dots$, luego se pasa al siguiente número que no haya sido tachado, que en este caso es $3$, y se repite el proceso.

Por ejemplo, si quisiéramos trabajar hasta $12$, al inicio tendríamos:

\begin{center}
\small
\begin{tabular}{|c|c|c|c|c|c|c|c|c|c|c|}
\hline
2 & 3 & 4 & 5 & 6 & 7 & 8 & 9 & 10 & 11 & 12 \\
\hline
\end{tabular}
\end{center}

Después de procesar el primo $2$, se tachan sus múltiplos:

\begin{center}
\small
\begin{tabular}{|c|c|c|c|c|c|c|c|c|c|c|}
\hline
2 & 3 & \st{4} & 5 & \st{6} & 7 & \st{8} & 9 & \st{10} & 11 & \st{12} \\
\hline
\end{tabular}
\end{center}

Luego el siguiente número no tachado es $3$, por lo que ahora se eliminan sus múltiplos que todavía sigan activos, como $9$. Este método funciona porque si un número no ha sido tachado por ningún primo más pequeño, entonces no puede tener divisores primos menores y por lo tanto debe ser primo.

Al terminar este proceso para los números hasta $12$, el estado final queda así:

\begin{center}
\small
\begin{tabular}{|c|c|c|c|c|c|c|c|c|c|c|}
\hline
2 & 3 & \st{4} & 5 & \st{6} & 7 & \st{8} & \st{9} & \st{10} & 11 & \st{12} \\
\hline
\end{tabular}
\end{center}

Los números que permanecen sin tachar son justamente los primos: $2,3,5,7,11$.

Además, no es necesario revisar todos los números hasta $n$. Basta continuar mientras $i^2 \le n$. La razón es la misma que en la prueba de primalidad: si un número compuesto tiene un divisor mayor que $\sqrt{n}$, entonces el otro divisor correspondiente ya es menor que $\sqrt{n}$ y debió haberse encontrado antes.

Tampoco hace falta empezar a tachar desde $2i$. Cuando llegamos al primo $i$, todos los múltiplos menores que $i^2$ ya fueron tachados por algún primo más pequeño. Por eso la implementación suele comenzar en $i^2$.

\begin{lstlisting}[language=C++]
vector<bool> prime(n + 1, true);
prime[0] = prime[1] = false;
for (int i = 2; i * i <= n; i++) {
    if (prime[i]) {
        for (int j = i * i; j <= n; j += i) {
            prime[j] = false;
        }
    }
}
\end{lstlisting}
 
La complejidad es $O(n \log \log n)$, prácticamente lineal. Funciona bien hasta $n \approx 10^7$ con memoria razonable, si se necesitan los primos hasta $10^8$ o el menor factor primo de cada número (útil para factorizar rápido), existe una criba lineal $O(n)$ que calcula el menor factor primo de cada entero. Se puede consultar en cp-algorithms.com.

\section{Teorema de Bézout}

\textbf{Propiedad.} Para cualesquiera enteros $a$ y $b$ no ambos cero, existen enteros $x$ e $y$ tales que
\[
ax + by = \gcd(a, b)
\]

Por ejemplo, si $a = 12$ y $b = 18$, entonces $\gcd(12,18) = 6$ y una posible elección es
\[
12(-1) + 18(1) = 6.
\]
Así, en este caso los coeficientes de Bézout pueden ser $x = -1$ y $y = 1$.
 
Los coeficientes $x$ e $y$ se llaman \textbf{coeficientes de Bézout} y no son únicos: si $(x_0, y_0)$ es una solución, entonces $(x_0 + k\frac{b}{\gcd}, \; y_0 - k\frac{a}{\gcd})$ también lo es para cualquier entero $k$.
 
\textbf{Propiedad.} La ecuación
\[
ax + by = c
\]
tiene solución entera si y sólo si $\gcd(a,b) \mid c$.

Por ejemplo, la ecuación
\[
12x + 18y = 6
\]
sí tiene solución entera porque $\gcd(12,18) = 6$ y claramente $6 \mid 6$. En cambio,
\[
12x + 18y = 5
\]
no tiene solución entera, ya que $6 \nmid 5$.
 
Los coeficientes de Bézout se calculan con el \textbf{algoritmo de Euclides extendido}. Mientras que el algoritmo de Euclides solo calcula $\gcd(a,m)$, su versión extendida además encuentra enteros $x, y$ tales que
\[
ax + my = \gcd(a,m)
\] 
La idea del algoritmo extendido es seguir los mismos pasos que Euclides pero rastrear cómo se expresa cada residuo como combinación lineal de $a$ y $m$. En cada paso de la división $r_{i-1} = q_i \cdot r_i + r_{i+1}$ se actualizan los coeficientes:
\[
x_{i+1} = x_{i-1} - q_i \cdot x_i \qquad y_{i+1} = y_{i-1} - q_i \cdot y_i
\]
partiendo de $(x_0, y_0) = (1, 0)$ para $a$ y $(x_1, y_1) = (0, 1)$ para $m$.

Veamos un ejemplo concreto con $a = 30$ y $m = 18$.

\textbf{Paso 1.} Aplicamos el algoritmo de Euclides:
\[
30 = 1 \cdot 18 + 12
\]
\[
18 = 1 \cdot 12 + 6
\]
\[
12 = 2 \cdot 6 + 0.
\]
Como el último residuo no nulo es $6$, concluimos que
\[
\gcd(30,18) = 6.
\]

\textbf{Paso 2.} Escribimos ese $6$ en función de los residuos anteriores:
\[
6 = 18 - 1 \cdot 12.
\]

\textbf{Paso 3.} Sustituimos ahora la expresión de $12$ que apareció en el primer paso:
\[
12 = 30 - 1 \cdot 18.
\]
Al reemplazarla en la ecuación anterior, obtenemos
\[
6 = 18 - (30 - 18).
\]

\textbf{Paso 4.} Simplificamos:
\[
6 = -30 + 2 \cdot 18.
\]
Por lo tanto, una elección válida de coeficientes de Bézout es
\[
x = -1, \qquad y = 2.
\]
 
\begin{lstlisting}[language=C++]
// calcula gcd(a,b) y encuentra x,y tales que ax + by = gcd(a,b)
long long gcdExt(long long a, long long b, long long &x, long long &y) {
    //Caso base pues gcd(a,0)=a
    if (b == 0) {
        x = 1; y = 0; 
        return a;
    }
    long long x1, y1;
    long long g = gcdExt(b, a % b, x1, y1);
    x = y1;
    y = x1 - (a / b) * y1;
    return g;
}
\end{lstlisting}
 
La complejidad es la misma que Euclides: $O(\log \min(a,b))$.

\section{Inverso multiplicativo modular}
 
Aunque pueda parecer que las operaciones con congruencias son muy flexibles, hay que tener cuidado, pues aún no hemos mencionado nada sobre la división. En aritmética ordinaria, el inverso multiplicativo de un número $a$ es $\frac{1}{a}$, es decir, el número tal que $a \cdot \frac{1}{a} = 1$. Por ejemplo, el inverso multiplicativo de $2$ es $\frac{1}{2}$, ya que
\[
2 \cdot \frac{1}{2} = 1.
\]
En aritmética modular necesitamos algo análogo: dado $a$ y un módulo $m$, queremos encontrar un entero $a^{-1}$ tal que
\[
a \cdot a^{-1} \equiv 1 \pmod{m}
\]
Por ejemplo, el inverso de $3$ módulo $7$ es $5$, pues
\[
3 \cdot 5 = 15 \equiv 1 \pmod{7}.
\]
Esto es útil en el ICPC cada vez que necesitamos ``dividir'' dentro de una operación modular. Como la división directa no está definida en modulo $m$, la reemplazamos por multiplicar por el inverso:
\[
\frac{a}{b} \bmod m \;=\; a \cdot b^{-1} \bmod m
\]
Por ejemplo, al calcular combinaciones
\[
\binom{n}{k} = \frac{n!}{k!(n-k)!}
\]
módulo un primo $p$, necesitamos dividir factoriales, y esa división se realiza multiplicando por inversos modulares. Si quisiéramos calcular, por ejemplo,
\[
\binom{5}{2} \bmod 7,
\]
no dividimos directamente entre $2!$ ni entre $3!$ dentro del módulo, sino que escribimos
\[
\binom{5}{2} \equiv 5! \cdot (2!)^{-1} \cdot (3!)^{-1} \pmod{7}.
\]
Este tipo de situaciones aparece con mucha frecuencia en ICPC, porque los números involucrados suelen crecer demasiado rápido y por eso es común trabajar módulo algún primo grande como $10^9+7$.
\\\\
Es importante mencionar que este inverso no siempre existe. Tomemos por ejemplo el caso de $m=4$ y $a=2$. Si existiera un inverso de $2$ módulo $4$, debería existir algún entero $b$ tal que
\[
2b \equiv 1 \pmod{4}.
\]
Pero si revisamos los posibles residuos módulo $4$, obtenemos:
\[
2 \cdot 0 \equiv 0,\qquad
2 \cdot 1 \equiv 2,\qquad
2 \cdot 2 \equiv 0,\qquad
2 \cdot 3 \equiv 2 \pmod{4}.
\]
Es decir, al multiplicar por $2$ módulo $4$ solo pueden aparecer los residuos $0$ y $2$, nunca el residuo $1$. Por eso en este caso no existe inverso modular.
\\\\El inverso de $a$ módulo $m$ existe si y sólo si $\gcd(a, m) = 1$, es decir, $a$ y $m$ son coprimos. Esto se puede probar facilmente, queremos $ax \equiv 1 \pmod{m}$, es decir, $ax - my = 1$ para algún entero $y$. Por el teorema de Bézout, esta ecuación tiene solución si y sólo si $\gcd(a,m) \mid 1$, lo que equivale a $\gcd(a,m) = 1$.
\\\\
Ahora, también podemos diseñar un algoritmo para encontrarle en caso de que exista, arriba vimos el algoritmo de euclides extendido, si lo aplicamos a $a$ y $m$, tenemos que como $\gcd(a,m) = 1$, esto nos da directamente $ax + my = 1$, es decir $ax \equiv 1 \pmod{m}$. Entonces $x$ es el inverso de $a$ módulo $m$. Entonces para implementar una función que nos de el inverso basta con llamar a \texttt{gcdExt}, es decir
\begin{lstlisting}[language=C++]
long long inv(long long a, long long m) {
    long long x, y;
    gcdExt(a, m, x, y);
    return (x % m + m) % m;  
\end{lstlisting}
Nótese que al final en lugar de devolver $x$, devolvemos $(x \% m + m) \% m$, esto es porque el coeficiente $x$ que devuelve \texttt{gcdExt} puede ser negativo. Por ejemplo, el inverso de $3$ módulo $7$ es $5$, pero el algoritmo puede devolver $x = -2$, que es equivalente pues $-2 \equiv 5 \pmod 7$. La expresión \texttt{(x \% m + m) \% m} convierte cualquier $x$ a un numero congruente a $x$ positivo en $[0, m)$.
\\
\\Cuando $m = p$ es primo, el pequeño teorema de Fermat dice que $a^{-1} \equiv a^{p-2} \pmod{p}$, entonces podemos calcular el inversocon exponenciación rápida en $O(\log p)$.
\\

\subsection*{Inversos para cada número módulo $m$}

En ocasiones queremos calcular el inverso modular para cada número $1,2,...,m-1$, si realizamos el algoritmo que describimos anteriormente nótese que esto tendría complejidad $O(m log m)$, si el módulo $m$ es primo podemos construir un algoritmo que tenga complejidad $O(m)$. Tomemos $i$ un número en el intervalo $[1,m-1]$, entonces
\[m \, mod \, i = m - \lfloor \frac{m}{i} \rfloor \cdot i\]
podemos sacar módulo $m$ de ambos lados y multiplicar por $i^{-1} \cdot (m \, mod \, i)^{-1}$ de donde el lado el lado izquierdo queda $i^{-1} \cdot (m \, mod \, i)^{-1} \cdot (m \, mod \, i)=i^{-1}$ y el derecho $- \lfloor \frac{m}{i} \rfloor \cdot i \cdot i^{-1} \cdot (m \, mod \, i)^{-1} mod \, m=- \lfloor \frac{m}{i} \rfloor \cdot (m \, mod \, i)^{-1} \, mod \, m$ de donde obtenemos la ecuación
\[i^{-1} = - \lfloor \frac{m}{i} \rfloor \cdot (m \, mod \, i)^{-1} \, mod \, m\]
 
de donde la implementación es muy corta

\begin{lstlisting}[language=C++]
vector<long long> inversos(int m, long long p) {
    vector<long long> inv(m);
    inv[1] = 1;
    for (int i = 2; i < m; i++)
        inv[i] = (p - (p / i) * inv[p % i] % p) % p;
    return inv;
}
\end{lstlisting}

\section{Pequeño teorema de Fermat}

Diremos que dos enteros son \textbf{primos relativos} o \textbf{coprimos} si su máximo común divisor es $1$.

\textbf{Propiedad.} Si $p$ es primo y $a$ es coprimo con $p$, es decir, si $\gcd(a, p) = 1$, entonces:
\[
a^{p-1} \equiv 1 \pmod{p}
\]
 
Multiplicando ambos lados por $a^{-1}$ se obtiene $a^{p-2} \equiv a^{-1} \pmod{p}$, lo que da un método directo para calcular el inverso modular cuando el módulo es primo.
 
Una generalización importante es el \textbf{teorema de Euler}: si $\gcd(a,m)=1$ entonces $a^{\phi(m)} \equiv 1 \pmod{m}$, donde $\phi(m)$ es la función de Euler (cantidad de enteros en $[1,m]$ coprimos con $m$). El pequeño teorema de Fermat es el caso particular $m = p$ primo, donde $\phi(p) = p-1$.

\section{Congruencias lineales y Teorema Chino del Residuo} 
\subsection*{Congruencia lineal}

Ahora, ya que sabemos calcular inversos, esto nos ayudará a solucionar congruencias lineales, es decir dados $a,\, b, \, n$ enteros, encontrar $x$ tal que\[ax \equiv b \,( mod \, m)\]
al igual que los inversos, estas congruencias no siempre tienen solución, en particular fijemonos en el $g=gcd(a,m)$ estas congruencias tendrán solución si y sólo si $g | b$, nótemos que además en caso de que $g=1$ tenemos que $a$ tiene inverso $módulo \, m$ entonces podemos multiplicar por el inverso en ambos lados y así encontramos que
\[x \equiv b a^{-1} \, mod \, m\]
¿Qué sucede en el caso que $g | b$ pero $g$ no es $1$? Nótemos por definición $g | a$ y $g | m$ y estamos suponiendo que $g | b$ entonces existen $a_g, m_g$ y $c_g$, tal que $a=ga_g$, $m=gm_g$ y $b=gb_g$ de dónde $ax-b=g(a_gx-b_g)$ de dónde $m$ es factor de $ax-b$ si y sólo si $m_g$ es factor de $a_gx-b_g$ de donde tenemos que $a_gx \equiv b_g \, mod \, m_g $, además el $mcd(a_g, m_g)=1$ por haber dividido entre $g$ entonces llegamos al caso de arriba que ya vimos como resolver con inversos.
\subsection*{Teorema chino del residuo}
Ahora veremos un resultado muy importante en teoría de números llamado \textbf{teorema chino del residuo},una pregunta natural al trabajar con congruencias es al trabajar con sistemas de congruencias, ¿qué sucede cuando tenemos dos congruencias diferentes?, ¿cómo encontrar una $x$ que satisfaga ambas de las ecuaciones siguientes para $a_1$, $a_2$, $n_1$, $n_2$ enteros? 
\[x \equiv a_1 \, mod \, n_1\]
\[x \equiv a_2 \, mod \, n_2\]
nótese que estas congruencias implican que 
\[x= a_1 + n_1b_1\]
\[x= a_2 + n_2b_2\]
es decir que
\[a_1+n_1b_1=a_2+n_2b_2\]
\[n_1(-b_1)+n_2b_2=a_1-a_2\]
denotemos $d=mcd(n_1,n_2)$, por definición, divide el lado izquierdo, entonces si $x$ existe, $d$ también divide el lado derecho, es decir $a_1-a_2$, ahora, ya sabemos que podemos encontrar $x'$ y $y'$ tal que $n_1x' + n_2y'=d$ con el algoritmo de euclides extendido, también podemos multiplicar ambos lados por $\frac{a_1-a_2}{d}$ y tenemos
\[n_1x'\frac{a_1-a_2}{d} + n_2y'\frac{a_1-a_2}{d}=a_1-a_2\]
de donde sabemos que $b_1=-\frac{a_1-a_2}{d}x'$ y $b_2=\frac{a_1-a_2}{d}y'$ y entonces podemos simplemente substituir $b_1$ para tener que
\[x=a_1+n_1\frac{a_2-a_1}{d}x'\]
esto nos da una solución al sistema y en realidad, no solamente una, sean $x_1$ y $x_2$ dos soluciones diferentes, ya que $x_1\equiv x_2 \, mod \, n_1$ y $x_1\equiv x_2 \, mod \, n_2$, tendremos que esto equivale a que $x_1 \equiv x_2 \, mod \, mcm(n_1,n_2)$ entonces todas los enteros congruentes a $x_1$ módulo $mcm(n_1,n_2)$ serán soluciones, lo que nos da una infinidad de soluciones. 
\\\\
Aunque esta implementación se ve directa, recordemos que siempre tenemos que considerar los posibles overflow que podemos tener, para evitar esto, es recomendable hacer las operaciones sacando módulo para que sean más pequeñas las cantidades.
\\\\
Ahora, este algoritmo no es exclusivo de 2 congruencias, en realidad si tuvieramos $k$ congruencias $x \equiv a_i \, mod \, n_i$ con $i=1,\cdots,k$  podemos ir aplicar el algoritmo a la primera y segunda y obtenemos una congruencia del tipo $x \equiv s \, mod \, mcm(n_1,n_2)$ y podemos unir esta con la tercera y así sucesivamente. Este algoritmo tendrá complejidad $O(klog(mcm(n_1,\cdots,n_k)))$.
\\\\
%EJERCICIO IMPLEMENTAR CRT PARA K CONGRUENCIAS
%https://codeforces.com/contest/982/problem/E

\section{Funciones multiplicativas}
 
Una función $f: \mathbb{Z}^+ \to \mathbb{Z}$ es \textbf{multiplicativa} si $f(ab) = f(a)f(b)$ siempre que $\gcd(a,b) = 1$. Esto permite calcularla sobre cualquier entero conociendo su valor en potencias de primos, usando la descomposición de primos del número.
 
\subsection*{Función de Euler $\phi$}
 
$\phi(n)$ cuenta los enteros en $[1,n]$ coprimos con $n$, recordemos que dos números son coprimos si su maximo común divisor es 1. Por ser función multiplicativa cumple la propiedad de que si $a$ y $b$ son coprimos entonces $\phi(ab)=\phi(a)\phi(b)$, ya que todo número $n$ se puede ver como un producto de primos $n=p_1^{k_1}p_2^{k_2}...p_m^{k_m}$ podemos obtener $\phi(n)$ de cualquier numero calculando $\phi$ para cada primo $p_i^{k_i}$ en su descomposición. Entonces ahora tenemos que preguntarnos cómo es el comportamiento de $\phi$ con los primos y sus potencias. Nótese que si $p$ es un primo por definición es coprimo a todos los números en $[1,p-1]$, entonces 
\[\phi(p)=p-1\]
Ahora si en lugar de $p$, tenemos $p^k$, notemos que hay $p^k/p$ números en $[1,p^k]$ divisibles por $p$, es decir no coprimos con $p^k$ de donde los podemos restarlos y obtener que
\[\phi(p^k)=p^k-p^{k_1}\]
Una propiedad importante a recordar es la \textbf{propiedad de la suma de divisores} que nos dice que si sumamos los $\phi(d)$ para todos los $d$ divisores de $n$ obtenemos $n$
\[\sum_{d \mid n} \phi(d) = n\]
Para la implementación podemos hacerla de forma directa empezando $n$ y cuando encontramos divisores quitar la cantidad de numeros que lo comparten en $[1,n]$, nótese que por ser un algoritmo de factorización, tiene complejidad $O(\sqrt{n})$.
\begin{lstlisting}[language=C++]
long long phi(long long n) {
    long long res = n;
    for (long long p = 2; p * p <= n; p++)
        if (n % p == 0) {
            while (n % p == 0){
                n /= p;
            }
            res -= res / p;
        }
    return res;
}
\end{lstlisting}
Ahora, como pasó con los inversos, habrá momentos en los que queramos calcular la función $\phi$ para todos los números de $1$ a $n$, si realizamos el algoritmo de arriba para cada uno tendría complejidad $O(n\sqrt{n})$, en realidad podemos mejorar mucho esta complejidad haciendo una pequeña modificación a nuestro algoritmo. En el algoritmo de arriba para cada factor primo modificamos un solo numero, en este caso, usaremos una idea similar a la de la criba, para cada primo actualizaremos todos los numeros en $[1,n]$ de los que sea factor, esto nos reduce la complejidad a $O(n log(log(n)))$.
\begin{lstlisting}[language=C++]
vector<long long> cribaPhi(int n) {
    vector<long long> phi(n + 1);
    //inciamos todos los resultados
    for(int i = 0; i <= n; i++){
        phi[i] = i;
    }
    for (int i = 2; i <= n; i++)
        //revisamos si i es primo
        if (phi[i] == i){
            //actualizamos todos los numeros de los que es factor i
            for (int j = i; j <= n; j += i){
                phi[j] -= phi[j] / i;
            }
        }
    return phi;
}
\end{lstlisting}
En una sección previa vimos el pequeño teorema de Fermat, su generalización cuando $m$ no es primo, se le llama \textbf{Teorema de Euler} y nos dice que si $a$ y $m$ son primos relativos entonces
\[a^{\phi(m)}\equiv 1 \,(mod \, m)\]
Nótese que de lo anterior tenemos que si $a$ y $m$ son primos relativos entonces
\[a^n\equiv a^{n\phi(m)}\, (mod\, m)\]
que podemos simplificar a
 \[a^n\equiv a^{n\,(mod \,\phi(m))}\, (mod\, m)\]
esto nos ayuda a calcular más rápido $a^n$ cuando $n$ es muy grande.

\subsection*{Función de Möbius $\mu$}
 
$\mu(n)$ se define como:
\[
\mu(n) = \begin{cases}
1  & \text{si } n = 1 \\
(-1)^k & \text{si } n \text{ es producto de } k \text{ primos distintos} \\
0  & \text{si } n \text{ tiene algún factor cuadrado}
\end{cases}
\]
 
Su uso principal es la \textbf{inversión de Möbius} que nos dice que si $g(n) = \sum_{d \mid n} f(d)$, entonces podemos invertir las funciones de la siguiente forma
\[f(n) = \sum_{d \mid n} \mu(d)\, g(n/d)\]
Esto permite recuperar $f$ a partir de $g$ y viceversa. Ahora estudiemos qué sucede con la suma de $\mu$ sobre los divisores de un número como hicimos arriba con $\phi$.
\[f(n) = \sum_{d \mid n} \mu(d)\]
Notemos que $f(1)=1$ y $f(p^k)=0$ pues los divisores son potencias de $p$ o 1, cualquier potencia $p^r$ con $r>1$ tendrá que $\mu(p^r)=0$ por definición de la función, entonces sólo nos falta ver qué sucede con $f(p)$, como $\mu(p)=-1$ y $\mu(1)=1$, $f(p^k)=0$. De aquí que como $f(n)$ es multiplicativa se cumpla que

\[\sum_{d \mid n} \mu(d)
 = \begin{cases}
1  & \text{si } n = 1 \\
0  & \text{cualquier otro caso} 
\end{cases}
\]
Podemos simplificar aún más la definición con un poco de notación, $[P]$ será $1$ si $P$ es verdadero, y 0 en cualquier otro caso, por ejemplo $[gcd(4,8)=4]=1$, de forma que quedaría
\[\sum_{d \mid n} \mu(d)=[n=1]\]
esto nos ayudará a simplificar la notación en los siguientes ejemplos. Ahora, podría parecer que la función de Möbius es algo sin mucha aplicación, pero es una herramienta poderosa, hagamos un ejemplo. Nuestra tarea será encontrar la cantidad de pares $(x,y)$ con $gcd(x,y)=1$ en el intervalo $[1,n]$.
\\Lo primero que debemos hacer al encontrar un enunciado de este tipo es intentar expresar lo que se nos pide, notemos que este enunciado en realidad nos pide calcular
\[f(n)=\sum_{j=1}^{n} \sum_{i=1}^{n} [gcd(i,j)=1]\]
ingenuamente podríamos pensar en recorrer todos los pares, pero esto nos llevaría complejidad $O(n^2)$ que es una complejidad muy alta, veamos como hacerlo más eficiente, con la propiedad de arriba tenemos que $[gcd(i,j)=1]=\sum_{d \mid gcd(i,j)} \mu(d)$ y además en esta última suma es la suma sobre todos los $d$ que dividen a $gcd(i,j)$, que es lo mismo que hacer la suma sobre todos los $d$ en $[1,n]$ y sólo contar el resultado cuando $d$ divide a $gcd(i,j)$, es decir
\[\sum_{d \mid gcd(i,j)} \mu(d)=\sum_{d =1}^n [d|gcd(i,j)]\mu(d)\]
y recordemos que $d$ divide a $gcd(i,j)$ si y solo si $d$ divide a $i$ y $d$ divide a $j$, entonces tenemos
\[\sum_{d =1}^n [d|gcd(i,j)]\mu(d)=\sum_{d =1}^n [d|i][d|j]\mu(d)\]
sustituyendo tenemos que
\[f(n)=\sum_{j=1}^{n} \sum_{i=1}^{n} \sum_{d =1}^n [d|i][d|j]\mu(d)\]
ahora simplemente podemos reordenar la forma de sumar para obtener
\[\sum_{d =1}^n \mu(d)(\sum_{i =1}^n [d|i])(\sum_{j =1}^n [d|j])\]
y notemos que $\sum_{i =1}^n [d|i]=\sum_{j =1}^n [d|j]=\lfloor \frac{n}{d} \rfloor$ de donde tenemos
\[f(n)=\sum_{d =1}^n \mu(d)\lfloor \frac{n}{d} \rfloor^2\]
podemos usar una criba parecida a la que hicimos con $phi$ para tener calculado $\mu$ para todos los números en $[1,n]$ y simplemente hacer un ciclo de $1$ a $n$ calculando la suma, esto nos simplifica la complejidad a $O(n)$. Vamos a implementar la criba:

\begin{lstlisting}[language=C++]
vector<int> cribaMu(int n) {
    vector<int> mu(n + 1, 1);
    vector<bool> compuesto(n + 1, false);
    vector<int> primos;
    for (int i = 2; i <= n; i++) {
        //revisar si es primo
        if (!compuesto[i]) { 
            primos.push_back(i); mu[i] = -1;
        }
        for (int p : primos) {
            if (i * p > n){
                break;
            }
            compuesto[i * p] = true;
            //revisar si tiene cuadrados
            if (i % p == 0) {
                mu[i * p] = 0;
                break;
            }
            //mu(ip)=mu(i)mu(p)=-mu(i)
            mu[i * p] = -mu[i];
        }
    }
    return mu;
}
\end{lstlisting}
\section{Logaritmo discreto y Baby-step Giant-step}
El \textbf{logaritmo discreto} se refiere al problema de dados $a$, $b$ y $m$, cómo encontrar $x$ tal que $a^x \equiv b \pmod{m}$. Este logaritmo puede no existir, por ejemplo $2^x\equiv 3 (mod 7)$. Nótese que por fuerza bruta podemos intentar probar $x = 0, 1, \ldots, m-1$ en $O(m)$, pero esto se puede hacer más eficientemente. El algoritmo \textbf{Baby-step Giant-step (BSGS)} lo resuelve en $O(\sqrt{m})$.
\\Consideremos la siguiente ecuación con $gcd(a,m)=1$
\[a^x \equiv b \pmod{m}\]
Nótese que cualquier $x$ en $[1,m-1]$ se puede escribir como $x=i\lceil \sqrt{m} \rceil - j $ con $1\leq i,j \leq \lceil \sqrt{m} \rceil$, de aquí que queramos resolver
\[a^{i\lceil \sqrt{m} \rceil - j} \equiv b \pmod{m}\]
 ahora, como $gcd(a,m)=1$ tiene inverso, entonces podemos considerar 
 \[a^{i\lceil \sqrt{m} \rceil} \equiv ba^{j} \pmod{m}\]
Este algoritmo se llama baby-step giant-step porque cuando recorremos valores de $i$ el $x$ da un brinco muy gigante, en cambio al variar los valores de $j$ se da brincos muy pequeños, lo que haremos será precalcular y guardar en un mapa todos los valores de $b \cdot a^j \bmod m$ para $j = 0, 1, \ldots, \lceil\sqrt{m}\rceil$, usaremos exponenciación binaria. Después calcularemos lo mismo para cada $i = 1, 2, \ldots, \lceil\sqrt{m}\rceil$, calcular $a^{i\lceil\sqrt{m}\rceil} \bmod m$ y buscar si está en el mapa, es decir si hay solución.
 
\begin{lstlisting}[language=C++]
long long bsgs(long long a, long long b, long long m) {
    a %= m; b %= m;
    if (b == 1 || m == 1){
        return 0;
    }
    
    //calculamos solo una vez el techo de la raiz 
    long long n = ceil(sqrt((double)m));
    unordered_map<long long, long long> baby;
 
    //baby steps guardar b*a^j para j = 0..n
    long long aj = b;
    for (long long j = 0; j <= n; j++) {
        baby[aj] = j;
        aj = aj * a % m;
    }
 
    //Usamos la exponenciacion binaria que vimos antes
    long long an = power(a, n, m);
    long long cur = an;
    for (long long i = 1; i <= n; i++) {
        if (baby.count(cur)) {
            //recuperar el exponente 
            long long x = i * n - baby[cur];
            if (x > 0){
                return x;
            }
        }
        cur = cur * an % m;
    }
    return -1;
}
\end{lstlisting}
Este algoritmo también puede usarse para el caso general (módulo no primo o $\gcd(a,m) > 1$) existe una versión extendida que reduce el problema al caso coprimo ver \href{https://cp-algorithms.com/algebra/discrete-log.html#when-a-and-m-are-not-coprime}{caso general}.
 
\section{Ejercicios}
%https://codeforces.com/problemset/problem/1553/A
\textbf{[3.1]} Definamos $S(x)$ como la suma de digitos de un número $x$ en sistema decimal. Por ejemplo $S(11)=2$, $S(1002)=3$, $S(5)=5$.
\\Diremos que un entero $x$ es interesante si $S(X+1)\leq S(X)$. En cada test se te dará un número $n$. Tu tarea es calcular cuántos números $x$ cumplen que $1\leq x \leq n$ y $x$ es interesante.
\\\\
\textbf{Input:}
\\
La primera línea contiene un entero $t$ ($1 \leq t \leq 1000$) — número de casos.
\\
Después le siguen $t$ líneas, la $i$-ésima línea contiene un número entero $n$ ($1\leq n \leq 10^9$).
\\\\
\textbf{Output:}
\\
Imprima $t$ enteros, el $i$-ésimo entero deberá ser la respuesta para el $i$-ésimo caso.
\\\\
\textbf{Ejemplos:}
\begin{table}[h !]
\begin{tabular}{lllll}
\cline{1-2}
\multicolumn{1}{|l|}{\textbf{Entrada}}                                                        & \multicolumn{1}{l|}{\textbf{Salida}}                                                  &  &  &  \\ \cline{1-2}
\multicolumn{1}{|l|}{\begin{tabular}[c]{@{}l@{}}5\\ 1\\ 9\\ 10\\ 34\\ 880055535\end{tabular}} & \multicolumn{1}{l|}{\begin{tabular}[c]{@{}l@{}}0\\ 1\\ 1\\ 3\\ 88005553\end{tabular}} &  &  &  \\ \cline{1-2}
                                                                                              &                                                                                       &  &  &  \\
                                                                                              &                                                                                       &  &  & 
\end{tabular}
\end{table}
\\
%https://codeforces.com/contest/1916/problem/B
\textbf{[3.2]} Un número $1\leq x \leq 10^9$ es elegido. Te darán dos enteros $a$ y $b$, los dos divisores más grandes del número $x$ y también se cumple que $1\leq a < b < x$.
\\
Para los números dados $a$, $b$ tienes que encontrar el valor de $x$.
\\\\
\textbf{Input:}
\\
Cada test contiene varios casos. La primera línea contiene un entero $t$ ($1\leq t \leq 10^4$), el número de casos. Luego la descripción de los casos.
\\ La única línea de cada caso contiene dos enteros $a,b \, (1 \leq a < b \leq 10^9)$.
\\ Se garantiza que $a,b$ son los divisores más grandes para algún número $1\leq x\leq 10^9$.
\\\\
\textbf{Output:}
\\
Para cada caso de prueba, imprime el número $x$, tal que $a$ y $b$ son los mayores divisores del número $x$.
Si hay varias respuestas, imprime cualquiera de ellas.
\\\\
\textbf{Ejemplos:}
\begin{table}[h !]
\begin{tabular}{lllll}
\cline{1-2}
\multicolumn{1}{|l|}{\textbf{Entrada}}                                                                            & \multicolumn{1}{l|}{\textbf{Salida}}                                                            &  &  &  \\ \cline{1-2}
\multicolumn{1}{|l|}{\begin{tabular}[c]{@{}l@{}}8\\ 2 3\\ 1 2\\ 3 11\\ 1 5\\ 5 10\\ 4 6\\ 3 9\\ 250000000 500000000\end{tabular}} & \multicolumn{1}{l|}{\begin{tabular}[c]{@{}l@{}}6\\ 4\\ 33\\ 25\\ 20\\ 12\\ 27\\ 1000000000\end{tabular}} &  &  &  \\ \cline{1-2}
                                                                                                                  &                                                                                                 &  &  &  \\
                                                                                                                  &                                                                                                 &  &  & 
\end{tabular}
\end{table}

%https://codeforces.com/problemset/problem/2097/C
\textbf{[3.3]} Dado un entero $n$, calcula el $n$-ésimo número de Fibonacci módulo $10^9+7$. Recuerda que la sucesión de Fibonacci está dada por
\[
F_0 = 0,\qquad F_1 = 1,\qquad F_n = F_{n-1}+F_{n-2}\text{ para } n\ge 2.
\]
La idea de este ejercicio es resolverlo usando exponenciación rápida, en particular mediante la potencia de la matriz
\[
\begin{pmatrix}
1 & 1 \\
1 & 0
\end{pmatrix},
\]
ya que esto permite obtener la respuesta en tiempo $O(\log n)$.
\\\\
\textbf{Input:}
\\
La única línea contiene un entero $n$ ($0 \leq n \leq 10^{18}$).
\\\\
\textbf{Output:}
\\
Imprime el valor de $F_n \bmod (10^9+7)$.
\\\\
\textbf{Ejemplos:}
\begin{table}[h !]
\begin{tabular}{lllll}
\cline{1-2}
\multicolumn{1}{|l|}{\textbf{Entrada}} & \multicolumn{1}{l|}{\textbf{Salida}} &  &  &  \\ \cline{1-2}
\multicolumn{1}{|l|}{7}                & \multicolumn{1}{l|}{13}              &  &  &  \\ \cline{1-2}
\end{tabular}
\end{table}